\documentclass[12pt,a4paper]{article}
\usepackage{amsmath,amssymb,amsthm}
\usepackage{hyperref}
\usepackage{geometry}
\usepackage{enumitem}
\geometry{margin=2.5cm}

\newtheorem{theorem}{Theorem}[section]
\newtheorem{corollary}[theorem]{Corollary}
\theoremstyle{definition}
\newtheorem{definition}[theorem]{Definition}
\theoremstyle{remark}
\newtheorem{remark}[theorem]{Remark}

\title{Discrete--Continuous Relation in Recognition Science:\\
How Spacetime Emerges from a Discrete Ledger}

\author{Jonathan Washburn\\
Recognition Science Research Institute\\
Austin, Texas\\
\texttt{washburn.jonathan@gmail.com}}

\date{March 2026}

\begin{document}
\maketitle

\begin{abstract}
We derive the discrete--continuous relation in Recognition Science.
States live on $\mathbb{Z}^D$; ambient topology is $\mathbb{R}^D$.
$J(x)$ enforces discreteness: only discrete lattice configurations
are stable minima.  Continuum emerges in the large-$N$ limit.
$\gamma$ measures the discrete--continuum gap.  $D = 3$ is forced by
8-tick epoch, linking, and Clifford algebra.
\end{abstract}

\section{Introduction}
\label{sec:intro}

Is spacetime fundamentally discrete or continuous?  LQG postulates
discrete area/volume~\cite{Rovelli}; causal sets use a discrete
poset~\cite{Sorkin}; string theory requires smooth background.  RS
resolves: the \emph{ledger} is discrete (voxels, ticks, integer
entries); the \emph{ambient topology} is continuous.  Continuous
mathematics emerges in the large-$N$ limit.  $J(x) =
\tfrac{1}{2}(x + x^{-1}) - 1$ is defined on $\mathbb{R}_{>0}$, but
stable minima lie on the discrete $\varphi$-ladder.  $k_R = \ln\varphi$
is the minimum non-trivial bit cost.

\section{Lattice and Ambient Space}
\label{sec:lattice}

\begin{definition}[Lattice State]
$\mathrm{LatticeState}(D) := \mathrm{Fin}(D) \to \mathbb{Z}$.
\end{definition}

\begin{definition}[Ambient Space]
$\mathrm{AmbientSpace}(D) := \mathrm{Fin}(D) \to \mathbb{R}$.
\end{definition}

\begin{definition}[Embedding]
The lattice embeds naturally: $\mathbb{Z}^D \hookrightarrow \mathbb{R}^D$.
\end{definition}

\section{Emergence of the Continuum from the Discrete Ledger}
\label{sec:emergence}

The RS ledger records discrete entries (integer scale, 8-tick phase).
No continuous degrees of freedom exist fundamentally.

\begin{theorem}[Continuum as Limit]
As $N \to \infty$, discrete sums $\sum_{k=1}^N f(k)$ approximate
continuum integrals; the finite remainder is $\gamma$.
\end{theorem}

\begin{remark}
$J(x)$ is defined on $\mathbb{R}_{>0}$, but the ledger stores only
discrete configurations.  Continuum domain permits variational
analysis; discrete ledger is the physical substrate.
\end{remark}

\section{$\mathbb{Z}_8 \to U(1)$: Discrete Phases to Continuum Gauge}
\label{sec:z8-u1}

The 8-tick epoch ($2^D = 8$) restricts phases to $\theta_k = k\pi/4$.

\begin{theorem}[Discrete Phase Group]
8-tick phases form $\mathbb{Z}_8$; continuum limit (coarse-grained
time) yields $\mathbb{Z}_8 \to U(1)$.
\end{theorem}

\begin{corollary}
Electromagnetic $U(1)$ emerges from 8-tick structure.  Hypercharge is
the continuum limit of $\mathbb{Z}_8$ phase invariance.
\end{corollary}

\section{$J$-Cost and Discreteness}
\label{sec:j-cost}

\begin{theorem}[Continuous Perturbation]
For any $\varepsilon > 0$, there exists a continuous state $x$ with
$J(x) < \varepsilon$.
\end{theorem}

\begin{theorem}[Discrete Cost Gap]
$J(2) > 0$ and $J(1/2) > 0$: non-trivial discrete states have positive
$J$-cost.
\end{theorem}

\begin{theorem}[Resolution]
States are discrete ($\mathbb{Z}^D$); topology is continuous
($\mathbb{R}^D$).  Discrete lattice sites are the stable minima.
\end{theorem}

\section{Lattice to Smooth Spacetime}
\label{sec:lattice-spacetime}

\begin{theorem}[Continuum Geometry]
In the large-$N$ limit, discrete ledger curvature (second variation of
$J$-cost) tends to the Einstein tensor.  EFE emerge with
$\kappa = 8\varphi^5$~\cite{RSEFE}.
\end{theorem}

\begin{remark}
Gravity is emergent; no graviton.  Discrete lattice + 8-tick becomes
smooth $(3+1)$-dimensional spacetime.
\end{remark}

\section{The Euler--Mascheroni Constant as Discrete--Continuum Gap}
\label{sec:gamma}

\begin{definition}[Euler--Mascheroni Constant]
$\gamma = \lim_{n\to\infty}(H_n - \ln n) \approx 0.5772$,
$H_n = \sum_{k=1}^n 1/k$.
\end{definition}

\begin{theorem}[Discrete Excess]
$\gamma$ measures the gap between discrete $H_n$ (cost of $n$ events)
and continuum $\ln n$.  Irreducible cost of discreteness~\cite{RSEuler}.
\end{theorem}

\begin{corollary}
RS renormalization: $\sum_{k=1}^\Lambda 1/k - \ln\Lambda \to \gamma$.
\end{corollary}

\section{Comparison with Other Discrete Approaches}
\label{sec:comparison}

\textbf{LQG}: Discreteness at Planck scale.  RS differs: ledger discrete
at all scales; smooth geometry emerges via large-$N$ and $J$-cost
minimization.  \textbf{Causal sets}: Discrete poset; continuum via
sprinkling.  RS differs: explicit algebraic content (8-tick,
$\varphi$-ladder) determines \emph{which} continuum emerges.
\textbf{Lattice QCD}: Continuum by tuning.  RS: no tuning; discrete
structure forces the limit uniquely.

\section{Predictions and Falsifiers}
\label{sec:falsifiers}

(1) Continuous spacetime below $\lambda_{\mathrm{rec}}$ challenges RS.
(2) 8-tick harmonics in EEG/oscillation spectra; absence disfavors RS.
(3) $\gamma$-dependent corrections in running couplings; exclusion
falsifies RS.  (4) No fundamental graviton; direct detection
contradicts RS.

\section{$D = 3$ Forcing}
\label{sec:d3}

\begin{theorem}[Dimension Forcing]
$D = 3$ is the unique dimension satisfying: (i) $2^D = 8$ (8-tick
epoch); (ii) non-trivial linking (knots require $D = 3$); (iii)
$\mathrm{Cl}_3 \cong M_2(\mathbb{C})$ (2-component spinors); (iv)
$\mathrm{Spin}(3) \cong SU(2)$.
\end{theorem}

\section{Conclusion}
\label{sec:conclusion}

The ledger is discrete; topology is continuous; continuum emerges in
the large-$N$ limit.  $\gamma$ quantifies the gap.  $D = 3$ is forced.
Falsifiable (\S\ref{sec:falsifiers}).

\begin{thebibliography}{9}
\bibitem{Rovelli} C.~Rovelli, \emph{Quantum Gravity} (Cambridge Univ.\ Press, 2004).
\bibitem{Sorkin} R.~D. Sorkin, Causal sets: Discrete gravity, in \emph{Approaches to Quantum Gravity}, ed.\ D.~Oriti (Cambridge Univ.\ Press, 2009).
\bibitem{RSEFE} J.~Washburn, Einstein field equations from Recognition Science, \texttt{RS\_Einstein\_Field\_Equations.tex}.
\bibitem{RSEuler} J.~Washburn, The Euler--Mascheroni constant in Recognition Science, \texttt{RS\_Euler\_Mascheroni\_Role.tex}.
\end{thebibliography}

\end{document}
